\section{Seven-vertex separators}\label{sec:seven-cuts}

This section is computation-free.  A connected subgraph of $G-S$ is
\emph{full at $S$} if it is adjacent to every vertex of $S$.
We retain the exact equality partition of a boundary colouring, because
its mere number of colours does not suffice for gluing.

\begin{lemma}\label{lem:reflection}
Let $V(G)=L\mathbin{\dot\cup}S\mathbin{\dot\cup}R$, with $L,R$ nonempty
and anticomplete, and suppose every proper minor of $G$ is six-colourable.
Let $\Pi$ partition $S$ into independent blocks.  Suppose that all its
blocks, except some singleton blocks whose vertices form a clique, can
be assigned injectively to disjoint connected subgraphs of $G[L]$ full
at $S$, with at least one block assigned.  Then $G[R\cup S]$ has a
six-colouring with equality partition on $S$ exactly $\Pi$.
\end{lemma}

\begin{proof}
For each assigned block $B_j$ and its full subgraph $P_j$, contract the
connected set $B_j\cup V(P_j)$.  These sets are disjoint.  Their images,
together with the retained singleton blocks, form a clique: fullness
supplies all adjacencies involving an assigned block, and the retained
singletons form a clique by hypothesis.  At least one edge is contracted,
so the resulting minor is proper.  Six-colour it and pull back only to
$R\cup S$, assigning a block the colour of its image.  Each block is
independent and every edge to a retained vertex was represented in the
minor, so the pullback is proper.  The clique forces distinct colours on
distinct blocks, giving exactly $\Pi$.
\end{proof}

\begin{lemma}\label{lem:seven-small}
A graph on seven vertices with at least ten edges contains a $K_4^{-}$,
a house, or $K_{2,3}$ as a not necessarily induced subgraph.  A house is
a four-cycle with an additional vertex joined to the ends of one cycle
edge.
\end{lemma}

\begin{proof}
Suppose $Q$ contains none of the three.  Choose a maximum-degree vertex
$v$, put $d=d(v)$, $A=N(v)$, and $B=V(Q)\setminus N[v]$.
The graph $Q[A]$ is a matching, since a two-edge path there gives
$K_4^{-}$ with $v$.  Every vertex of $B$ has at most two neighbours in
$A$, since three would give $K_{2,3}$ with $v$.  For $d\geq5$,
\[
 |E(Q)|\leq d+\lfloor d/2\rfloor+2(6-d)+\binom{6-d}{2}\leq9.
\]
For $d\leq2$ there are at most seven edges.  Write
$\alpha=|E(Q[A])|$, $x=|E_Q(A,B)|$, $\beta=|E(Q[B])|$.

If $d=4$, then $\alpha\leq2$, $x\leq4$, $\beta\leq1$.
For $\alpha=2$, ten edges force a vertex of $B$ with two neighbours in
$A$.  If these are matched, there is a $K_4^{-}$; if they are on distinct
matching edges, there is a house.  For $\alpha\leq1$, ten edges force
$(\alpha,x,\beta)=(1,4,1)$.  Let $ab$ be the edge in $A$ and $c,t$ its
isolated vertices.  A two-set of neighbours in $A$ using $a$ or $b$
gives $K_4^{-}$ or a house.  Thus the two vertices $p,q$ of $B$ both see
$c,t$; their edge gives a $K_4^{-}$ on $p,q,c,t$.

If $d=3$, maximum degree gives exactly ten edges, so
$\alpha+x+\beta=7$ and $\alpha\leq1$.  For $\alpha=1$, no vertex of
$B$ can have two neighbours in $A$, by the same house or $K_4^{-}$
argument.  Hence $x=\beta=3$: $B$ is a triangle and its three edges to
$A$ have distinct ends, since a repetition gives $K_4^{-}$.  The edge
of $A$, its two matched vertices in $B$, and $v$ form a house.

It remains that $\alpha=0$.  If $\beta=3$, each vertex of $B$ already has two neighbours in $B$,
so $x\leq3$, a contradiction.  If $\beta=2$, write $B=pqr$ as a path.
Then $x=5$, and its cross-degrees are $2,1,2$.  Equal two-element
$A$-neighbourhoods of $p,r$ give $K_{2,3}$.  If they differ, the neighbour
of $q$ is either their common vertex, giving $K_4^{-}$, or one of the
two other vertices, giving a house.  If $\beta=1$, every vertex of $B$
has two neighbours in $A$.  Repeated two-sets give $K_4^{-}$ or
$K_{2,3}$.  Otherwise label the edge of $B$ as $pq$ and write
\[
 N_A(p)=\{a,b\},\quad N_A(q)=\{a,c\},\quad N_A(r)=\{b,c\}.
\]
Then $a,v,b,p$ is a four-cycle with roof $q$.  Finally $\beta=0$ would
force $x=7>6$.  Thus $Q$ has at most nine edges.
\end{proof}

\begin{lemma}\label{lem:three-cut-boundary}
Let $G$ satisfy the hypotheses of Theorem~\ref{thm:main}, and let $S$
be a seven-vertex cut.
Then $G-S$ has two or three components.  If it has three, $G[S]$ has a
proper three-colouring and every such colouring has class sizes $3,2,2$.
\end{lemma}

\begin{proof}
Every component is full at $S$ by seven-connectivity.  If there are at
least five, contract five of them to $c_1,\ldots,c_5$, label
$S=\{s_1,\ldots,s_7\}$, and use
\[
             \{c_i,s_i\}\ (1\leq i\leq5),\qquad\{s_6\},\{s_7\}.
\]
All but possibly the last pair are adjacent, giving $K_7^{-}$.

Suppose there are four components, contracted to $c_1,\ldots,c_4$.
If $G[S]$ has a path $xyz$, choose three further vertices $u,v,w$ of
$S$ and use
\[
 \{c_1\},\quad\{c_2,u\},\quad\{c_3,v\},\quad\{c_4,w\},
 \quad\{x\},\quad\{y\},\quad\{z\}.
\]
Only $\{x\}:\{z\}$ may be nonadjacent, again impossible.  Hence $G[S]$
is a matching and isolated vertices.  It has an isolated vertex $q$ and
a partition $S=I_1\mathbin{\dot\cup}I_2\mathbin{\dot\cup}\{q\}$ into
independent sets of orders $3,3,1$.  Split the four components into two
pairs, whose unions are $L,R$.  In each union the two full components
can be assigned to $I_1,I_2$ in Lemma~\ref{lem:reflection}, retaining
$q$.  The two colourings have the same exact boundary partition;
permuting colour names makes them agree on $S$ and yields a six-colouring
of $G$, a contradiction.  Thus there are at most three components.

Suppose there are three, with contracted vertices $c_1,c_2,c_3$.
If $G[S]$ had at least ten edges, Lemma~\ref{lem:seven-small} would
supply one of its three subgraphs.  For a $K_4^{-}$ subgraph, retain its
four vertices and merge each $c_i$ with one of the three unused boundary
vertices.  For a house or $K_{2,3}$, contract respectively the square
edge opposite the roof or any edge to obtain a four-bag $K_4^{-}$ model
using five boundary vertices.  Retain $c_1$ and merge $c_2,c_3$ with the
two unused boundary vertices.  In either case the seven sets form a
$K_7^{-}$ model.  Consequently $|E(G[S])|\leq9$.  The boundary also has
no $K_4$ subgraph, by the same construction using four singleton roots.

Every $K_4$-free graph on seven vertices with at most nine edges is
three-colourable.  Otherwise take a four-critical subgraph $F$; it has
minimum degree at least three.  On four vertices it is $K_4$.  On five
vertices parity gives a universal vertex, whose deletion leaves a
three-chromatic four-vertex graph, containing a triangle; this also gives
$K_4$.  On six vertices a degree at least four forces ten edges.  Otherwise
$F$ is cubic, with complement $C_6$ or $2C_3$; pairing consecutive
vertices in the former or taking the triangles in the latter gives a
three-colouring.  On seven vertices minimum degree three gives at least
eleven edges.  These contradictions prove the assertion.

If $G[S]$ were two-colourable, fix its partition into at most two
independent blocks.  For each component use the other two full components
in Lemma~\ref{lem:reflection} to reproduce the partition on that closed
component.  The three colourings align and glue, a contradiction.
Likewise, if any proper three-colouring has a singleton class, retain that
singleton and assign the other two blocks to the other two components.
The same gluing contradiction results.  Every class of a proper
three-colouring therefore has order at least two; their orders are
$3,2,2$.
\end{proof}

\begin{theorem}\label{thm:three-component-colouring}
Let $G$ be seven-connected, with every proper minor six-colourable.  Let
$S$ be a seven-vertex cut with exactly three components $D_1,D_2,D_3$.
If $G[S]$ has a proper colouring with class sizes $3,2,2$, then $G$ is
six-colourable.  No excluded-minor hypothesis is needed.
\end{theorem}

\begin{proof}
Fix the independent boundary blocks
$S=T\mathbin{\dot\cup}A\mathbin{\dot\cup}B$, with
$|T|=3$ and $|A|=|B|=2$.  Every $D_i$ is full at $S$.
Say that a component realises $(A,B)$ if it contains disjoint connected
subgraphs $X_A,X_B$, the first adjacent to both vertices of $A$, the
second to both vertices of $B$.  They can also be made adjacent: join
them by a shortest path in the component, split it at an edge, and absorb
the two halves into the respective subgraphs.

If two components realise $(A,B)$, then for every retained component
$D_i$ we may choose a different realising component $R$ and let $F$ be
the third.  Contract the three disjoint connected sets
\[
                X_A\cup A,\qquad X_B\cup B,\qquad F\cup T.
\]
Their images form a triangle, using the adjacency of $X_A$ and $X_B$
and fullness of $F$.  A six-colouring of this proper minor pulls back on
$G[D_i\cup S]$ with exact boundary partition $T|A|B$.  The three
colourings then align and glue.  Thus we may assume $D_1,D_2$ do not
realise $(A,B)$.

Fix one of these two components $D_i$.  Add four artificial terminals
$a_1^*,b_1^*,a_2^*,b_2^*$, each with the neighbours in $D_i$ of its
corresponding literal root.  There are no two disjoint paths joining the
$a$-pair and the $b$-pair, since their interiors would realise $(A,B)$.
Add the cycle edges in the displayed terminal order.  These do not
create such a crossing: a path using an edge between consecutive
terminals would meet a terminal of the other prescribed path.  Extend
on the same vertex set to an edge-maximal graph with no such crossing.
Humeau and Pous~\cite[Theorem~1.3]{HumeauPous2025} identify this completion
as a web with the four terminals as frame.  In this form a web consists
of a planar graph with the frame on a face and with cliques attached
behind facial triangles.

No clique inserted behind such a triangle can contain a vertex of $D_i$
outside the triangle.  Otherwise those original vertices have no
auxiliary neighbours outside the triangle.  Replace any artificial
terminal of the triangle by its literal root and add the three omitted
vertices of $T$.  These at most six actual vertices then separate a
nonempty part of $D_i$ from another component of $G-S$, contrary to
seven-connectivity.  Thus every original vertex of $D_i$ belongs to the
planar part of the web.  Replacing the artificial terminals by their
literal roots shows that
\[
                   H_i^+=G[D_i\cup A\cup B]+E(K_{A,B})
\]
is planar, with the added $K_{2,2}$ as an induced facial four-cycle.
The cycle is induced because $A$ and $B$ are independent.  These added
edges serve only in the planar extension; they are not asserted to be
edges of $G$ or used in a minor of $G$.

Contract in $G$ the two disjoint connected sets
\[
                        Q_0=D_1\cup T,\qquad Q_1=D_2\cup A.
\]
Their images are adjacent and each sees both retained vertices of $B$.
Six-colour this proper minor.  Pulling back on $G[D_3\cup S]$ gives a
colouring $c$ in which $T$ has colour $0$, $A$ a different colour $1$,
and both vertices of $B$ avoid $0,1$.

For $i=1,2$, precolour the induced four-cycle of $H_i^+$ by
$c|_{A\cup B}$.  It uses at most three colours from the five-colour
palette excluding $0$.  Diwan~\cite[Corollary~1]{Diwan2023} proves that
a proper precolouring of an induced cycle of length at most $2k-5$
using at most $k-1$ colours extends to a $k$-colouring of every planar
supergraph containing it.  With $k=5$ it extends to each $H_i^+$.
Use these extensions on $D_1,D_2$ and retain $c$ on $D_3\cup S$.
They agree on $A\cup B$ and avoid $0$ on $D_1,D_2$, so every edge to
$T$ is proper.  Distinct components are anticomplete.  The combined
colouring is a six-colouring of $G$.
\end{proof}
